\documentclass[12pt]{article}
\usepackage[T1]{fontenc}
\usepackage[utf8]{inputenc}
\usepackage{lmodern}
\usepackage{textcomp}
\usepackage{enumitem}
\usepackage{geometry}
\geometry{margin=1in}
\begin{document}
\begin{center}
Answer Key - Constitutional Law - Undergraduate Level Assessment 15 \
\small (Answers are ordered exactly as in the question paper.)
\end{center}

\section*{Multiple Choice}
\begin{enumerate}[label=\arabic*.]
\item \textbf{Answer:} C
\\ \textit{Options:} A) The purpose of immunity is to protect foreign Heads of State from embarrassment; B) Immunity protects a State from being invaded by another; C) Immunity shields States from being sued in the courts of other States; D) The purpose of immunity is to offer impunity in respect of all crimes
\\ \textit{Note:} The correct answer highlights that sovereign immunity is a legal doctrine that protects a state from being sued in the courts of other states, thereby upholding the principle of state sovereignty and respect among nations.
\item \textbf{Answer:} D
\\ \textit{Options:} A) Authority.; B) Charisma.; C) Co-operation.; D) Capitalism.
\\ \textit{Note:} Weber argues that the development of formally rational law is closely linked to the rise of capitalism, as it emphasizes efficiency, predictability, and calculability in legal systems that support economic transactions and the growth of capitalist society.
\item \textbf{Answer:} C
\\ \textit{Options:} A) Golder v UK (1978); B) A v UK (2009); C) Goodwin v UK (2002); D) Osman v UK (1998)
\\ \textit{Note:} The case of Goodwin v UK (2002) exemplifies the 'living instrument principle' as it demonstrates the European Court of Human Rights' approach to interpreting the European Convention on Human Rights in a progressive manner to adapt to contemporary societal changes, particularly in recognizing the rights of transgender individuals.
\item \textbf{Answer:} A
\\ \textit{Options:} A) Article 3; B) Article 8; C) Article 9; D) Article 11
\\ \textit{Note:} Article 3 of the European Convention on Human Rights, which prohibits torture and inhuman or degrading treatment, is considered an absolute right and thus is not subject to any limitations or qualifications, distinguishing it from qualified rights that can be restricted under certain circumstances.
\item \textbf{Answer:} B
\\ \textit{Options:} A) The Council of Europe was established in 1950 and consists of 27 member states; B) The Council of Europe was established in 1949 and consists of 47 member states; C) The Council of Europe was established in 1959 and consists of 34 member states; D) The Council of Europe was established in 1984 and consists of 19 member states
\\ \textit{Note:} The correct answer is accurate because the Council of Europe was indeed established in 1949 and currently comprises 47 member states, making it a significant entity in promoting human rights, democracy, and the rule of law in Europe.
\end{enumerate}

\section*{True / False}
\begin{enumerate}[label=\arabic*.]
\setcounter{enumi}{5}
\item \textbf{Answer:} False
\\ \textit{Options:} A) True; B) False
\\ \textit{Note:} An 'armed attack' must be a clear and substantial act of aggression that meets specific criteria under international law, and does not justify preemptive invasion based solely on potential threats.
\item \textbf{Answer:} False
\\ \textit{Options:} A) True; B) False
\\ \textit{Note:} Critical legal theorists argue that law is not a stable and determinate set of rules, but rather a fluid and subjective construct influenced by social, political, and economic factors.
\item \textbf{Answer:} True
\\ \textit{Options:} A) True; B) False
\\ \textit{Note:} Normative legal theory is concerned with how laws ought to be and the moral principles and political values that inform legal standards and judgments.
\item \textbf{Answer:} True
\\ \textit{Options:} A) True; B) False
\\ \textit{Note:} Leopold Pospisil identifies 'sanction' as a crucial component of law because it refers to the enforcement mechanisms that ensure compliance with legal norms, thereby playing a vital role in the functioning of any legal system.
\item \textbf{Answer:} True
\\ \textit{Options:} A) True; B) False
\\ \textit{Note:} The statement is true because Cicero believed that natural law is an inherent moral order that can be understood through human reason, aligning legal principles with the natural world.
\end{enumerate}

\section*{Long Form Response}
\begin{enumerate}[label=\arabic*.]
\setcounter{enumi}{10}
\item \textbf{Answer:} Catharine MacKinnon argues that the power dynamics between men and women are rooted in systemic structures that perpetuate male dominance and female subordination. She contends that true equality cannot be achieved merely through legal reforms; rather, it requires significant political pressure to challenge and change these entrenched power relations. Believing that political pressure can lead to genuine equality implies a recognition of the necessity for active resistance against patriarchal systems, which often resist change. However, this perspective also raises questions about the effectiveness of such pressure, as it may lead to token reforms rather than substantive transformation. Ultimately, while political pressure is essential for promoting women's rights, it must be coupled with a deeper societal commitment to dismantling the underlying inequalities that perpetuate gender disparities.
\\ \textit{Note:} correct because it accurately reflects MacKinnon's view that systemic power dynamics require active political resistance to address entrenched gender inequalities, rather than relying solely on legal reforms for achieving true equality.
\item \textbf{Answer:} The principle of sovereign equality among states is a foundational concept in international law, enshrined in Article 2, paragraph 1 of the UN Charter, which asserts that all states are equal as sovereign entities. This principle implies that regardless of a state's size, power, or economic status, each possesses the same legal rights and obligations under international law. In contemporary global relations, sovereign equality fosters mutual respect and cooperation, enabling smaller nations to participate meaningfully in international forums alongside more powerful states. However, it also raises challenges, as disparities in resources and influence can lead to unequal power dynamics in practice, complicating the ideal of genuine equality. Overall, the principle remains crucial in promoting a just and balanced international order.
\\ \textit{Note:} correct because it accurately identifies sovereign equality as a fundamental principle of international law, rooted in the UN Charter, and highlights its importance in promoting equal legal rights and fostering cooperation among states, regardless of their relative power or status.
\end{enumerate}

\vfill
\noindent\textit{End of Answer Key}
\end{document}